\section{Historias de usuario: Sprint 3}

Las tareas se identifican con los commits de la fuente del sprint. Las horas son estimaciones razonables para el cambio descrito; la fecha de inicio y fin corresponde al día del commit. Los commits de Claude se asignan a Rolando Sebastián Miranda por su integración de los PR 17 y 18 y el merge indicado.

\subsection{HU-12: Registrar la solicitud con un único certificado SIHA y revisarlo antes de enviarlo}
\historia{HU-12}{Operador}{Registrar la solicitud con un único certificado SIHA y revisarlo antes de enviarlo}{Alta}{Alta}{8 SP}{4 al 13 de noviembre}{Alex Jonathan Chisag; Jair León; Rolando Sebastián Miranda}{Como operador, quiero registrar una solicitud con un certificado SIHA único y revisar los datos leídos antes de enviarla, para corregir problemas antes de que comience el análisis.}{\item Dado un certificado SIHA válido, cuando se registra, entonces queda una solicitud visible con sus datos y estado.\item Dado un certificado con campos incorrectos, cuando se revisa, entonces aparece una lista de verificación por sección y no se permite registrarlo.\item Dado un certificado ya registrado, cuando se intenta enviarlo otra vez, entonces se informa de la solicitud existente.}{El flujo utiliza el formato SIHA-CH-01 y evita duplicar el análisis.}{Registro y revisión del certificado}
\begin{figure}[H]\centering\begin{tikzpicture}[actor/.style={draw,ellipse},caso/.style={draw,ellipse,minimum width=3.3cm},node distance=1.1cm and 2cm]\node[actor] (a) {Operador};\node[caso,right=of a,yshift=1cm] (sube) {Subir certificado};\node[caso,below=of sube] (revisa) {Revisar secciones};\node[caso,below=of revisa] (registra) {Registrar solicitud};\draw (a)--(sube);\draw (a)--(revisa);\draw (a)--(registra);\draw[dashed] (revisa)--node[right]{«incluye»}(registra);\end{tikzpicture}\caption{Casos de uso de HU-12}\end{figure}
\textbf{Desglose de tareas}
\tarea{T12.1}{HU-12}{Registrar solicitud con certificado único}{D}{6 h}{4 nov. 2026}{4 nov. 2026}{Rolando Sebastián Miranda}{Commit 55d6ac3: se conectó el certificado único con el flujo de homologación automática por contenidos.}
\tarea{T12.2}{HU-12}{Validar certificado antes del registro}{D}{5 h}{13 nov. 2026}{13 nov. 2026}{Alex Jonathan Chisag}{Commit f1b71d6: se añadió la revisión previa del certificado antes de guardarlo.}
\tarea{T12.3}{HU-12}{Sugerir materia oficial de destino}{M}{3 h}{13 nov. 2026}{13 nov. 2026}{Alex Jonathan Chisag}{Commit b1fa62d: se sugirió el nombre oficial de la materia de destino.}
\tarea{T12.4}{HU-12}{Guardar contexto certificado y matrícula}{D}{4 h}{4 nov. 2026}{4 nov. 2026}{Jair León}{Commit 4c8a20e: se guardó el contexto de materias certificadas y matrícula.}
\tarea{T12.5}{HU-12}{Mejorar lectura OCR y revisión del operador}{C}{4 h}{4 nov. 2026}{4 nov. 2026}{Jair León}{Commit 9655d53: se reforzó la lectura OCR y el flujo para que el operador revise los datos.}
\tarea{T12.6}{HU-12}{Listar materias de la carrera elegida}{C}{3 h}{4 nov. 2026}{4 nov. 2026}{Jair León}{Commit ae66e47: se corrigió la lista de materias de origen según la carrera seleccionada.}
\tarea{T12.7}{HU-12}{Ajustar errores de integración del certificado}{C}{3 h}{4 nov. 2026}{4 nov. 2026}{Rolando Sebastián Miranda}{Commit 4256e9f: se corrigió el conteo de migraciones de DB B para que las pruebas de integración encontraran la estructura.}
\tarea{T12.8}{HU-12}{Preparar certificados de prueba}{O}{4 h}{4 nov. 2026}{4 nov. 2026}{Rolando Sebastián Miranda}{Commit a793d91: se prepararon 26 certificados SIHA de prueba, un generador y sus resultados.}
\tarea{T12.9}{HU-12}{Aclarar avisos y campos del certificado}{M}{3 h}{4 nov. 2026}{4 nov. 2026}{Rolando Sebastián Miranda}{Commit 5a8845d: se ajustaron notas con coma, el aviso de sílabo ausente, el formato anterior y la carga incremental de datos.}
\tarea{T12.10}{HU-12}{Mostrar menú y gestión de contraseñas}{M}{4 h}{22 nov. 2026}{22 nov. 2026}{Alex Jonathan Chisag}{Commit 0a3f102: se ordenó el menú lateral, los avisos y la gestión de contraseñas; también se corrigieron mensajes de errores del flujo.}
\tarea{T12.11}{HU-12}{Marcar documentos como archivos binarios}{C}{2 h}{13 nov. 2026}{13 nov. 2026}{Alex Jonathan Chisag}{Commit 3b863e5: se indicó a Git que los PDF y las imágenes deben tratarse como archivos binarios.}
\tarea{T12.12}{HU-12}{Preparar docentes de ejemplo}{O}{2 h}{24 nov. 2026}{24 nov. 2026}{Rolando Sebastián Miranda}{Commit e5fcff3: se agregaron docentes de ejemplo para poder probar el flujo de revisión.}
\textbf{Lo que queda funcionando}
\begin{itemize}\item El operador revisa los datos del certificado y la lista de problemas antes de registrar.\item Las solicitudes muestran sus materias y estado; el certificado duplicado se relaciona con el registro existente.\item La lectura identifica las materias de la carrera elegida y conserva su contexto.\end{itemize}
\textbf{Pruebas de aceptación}
\begin{table}[H]\centering\small\begin{tabularx}{\textwidth}{|l|X|X|X|}\hline\rowcolor{grisTarjeta}\textbf{PA} & \textbf{Prueba} & \textbf{Entrada} & \textbf{Resultado esperado}\\\hline PA-12.1 & Solicitud registrada & Certificado SIHA válido & La solicitud aparece con sus datos y estado.\\\hline PA-12.2 & Lista de verificación & Certificado con datos que incumplen reglas & Se muestran las secciones y motivos; no se registra.\\\hline PA-12.3 & Certificado repetido & Mismo certificado cargado otra vez & Se avisa y se ofrece la solicitud existente.\\\hline\end{tabularx}\caption{Pruebas de aceptación de HU-12}\end{table}

\subsection{HU-13: Analizar con los sílabos oficiales de la UTA y explicar cada materia}
\historia{HU-13}{Docente}{Analizar con sílabos oficiales de la UTA y explicar cada materia}{Alta}{Alta}{8 SP}{4 al 24 de noviembre}{Rolando Sebastián Miranda; Alex Jonathan Chisag}{Como docente, quiero que el agente compare cada materia con los sílabos oficiales y explique sus resultados, para revisar una recomendación comprensible y actualizada.}{\item Dada una solicitud con materias de la UTA, cuando termina el análisis, entonces se explica el resultado de cada materia con sus datos y regla aplicada.\item Dado que el agente aún se prepara o no estaba disponible, cuando el servicio vuelve, entonces la solicitud se analiza automáticamente.}{El agente recomienda y los docentes deciden, según la regla vigente.}{Análisis por materia}
\begin{figure}[H]\centering\begin{tikzpicture}[actor/.style={draw,ellipse},caso/.style={draw,ellipse,minimum width=3.4cm},node distance=1.1cm and 2cm]\node[actor](a){Docente};\node[caso,right=of a,yshift=0.6cm](ver){Ver análisis};\node[caso,below=of ver](detalle){Revisar explicación};\node[caso,below=of detalle](reintento){Reintentar análisis};\draw(a)--(ver);\draw(a)--(detalle);\draw(a)--(reintento);\draw[dashed](ver)--node[right]{«incluye»}(detalle);\end{tikzpicture}\caption{Casos de uso de HU-13}\end{figure}
\textbf{Desglose de tareas}
\tarea{T13.1}{HU-13}{Cargar sílabos oficiales y explicar análisis}{D}{6 h}{4 nov. 2026}{4 nov. 2026}{Rolando Sebastián Miranda}{Commit 687f6bf: se conectó el agente a los sílabos oficiales UTA y se presentó el análisis explicado y el historial del certificado.}
\tarea{T13.2}{HU-13}{Simplificar la solicitud y asignar revisores}{M}{6 h}{4 nov. 2026}{4 nov. 2026}{Rolando Sebastián Miranda}{Commit c8abebd: se incorporaron los dos revisores docentes, el análisis al subir el PDF y una aplicación más sencilla.}
\tarea{T13.3}{HU-13}{Reintentar el análisis al estar listo el agente}{D}{4 h}{13 nov. 2026}{13 nov. 2026}{Alex Jonathan Chisag}{Commit fc4cae5: se agregó el reintento automático mientras el agente se preparaba.}
\tarea{T13.4}{HU-13}{Volver a analizar con reglas vigentes}{C}{4 h}{24 nov. 2026}{24 nov. 2026}{Rolando Sebastián Miranda}{Commit 28c495d: se habilitó el nuevo análisis con reglas vigentes y se liberaron certificados del formato anterior.}
\tarea{T13.5}{HU-13}{Estabilizar pruebas de calidad e integración}{C}{4 h}{13 nov. 2026}{13 nov. 2026}{Alex Jonathan Chisag}{Commit 9602631: se corrigió el redondeo decimal y se estabilizaron pruebas del bus y ESLint.}
\tarea{T13.6}{HU-13}{Ajustar contrato y ejecución de integración}{C}{2 h}{14 nov. 2026}{14 nov. 2026}{Alex Jonathan Chisag}{Commit 1d359be: se ajustaron pruebas de integración al evento v1.2 y a .NET 8.}
\textbf{Lo que queda funcionando}
\begin{itemize}\item El análisis presenta la explicación y recomendación por materia con información oficial disponible.\item Si el agente no estaba listo, la solicitud puede continuar al volver el servicio.\item La decisión queda para los docentes de acuerdo con la regla vigente.\end{itemize}
\textbf{Pruebas de aceptación}
\begin{table}[H]\centering\small\begin{tabularx}{\textwidth}{|l|X|X|X|}\hline\rowcolor{grisTarjeta}\textbf{PA} & \textbf{Prueba} & \textbf{Entrada} & \textbf{Resultado esperado}\\\hline PA-13.1 & Análisis explicado & Solicitud con materias UTA & Muestra explicación, regla y recomendación por materia.\\\hline PA-13.2 & Reintento automático & Agente vuelve a estar disponible & La solicitud continúa al análisis sin un nuevo registro.\\\hline\end{tabularx}\caption{Pruebas de aceptación de HU-13}\end{table}

\subsection{HU-14: Movilidad externa: comparar el sílabo de otra universidad}
\historia{HU-14}{Operador}{Comparar el sílabo de otra universidad}{Media}{Alta}{5 SP}{13 de noviembre}{Alex Jonathan Chisag}{Como operador, quiero adjuntar el sílabo de una materia aprobada en otra universidad, para que el agente pueda compararlo con el sílabo de destino de la UTA.}{\item Dada una solicitud de movilidad externa con sílabo válido, cuando se analiza, entonces el agente usa sus unidades y temas para comparar.\item Dada una materia sin sílabo externo, cuando se analiza, entonces se indica que no hay sílabo para comparar.}{El sílabo externo es opcional.}{Movilidad externa}
\begin{figure}[H]\centering\begin{tikzpicture}[actor/.style={draw,ellipse},caso/.style={draw,ellipse,minimum width=3.4cm},node distance=1.1cm and 2cm]\node[actor](a){Operador};\node[caso,right=of a,yshift=.6cm](adj){Adjuntar sílabo};\node[caso,below=of adj](anal){Analizar materia externa};\node[caso,below=of anal](ver){Ver comparación};\draw(a)--(adj);\draw(a)--(anal);\draw(a)--(ver);\draw[dashed](anal)--node[right]{«incluye»}(ver);\end{tikzpicture}\caption{Casos de uso de HU-14}\end{figure}
\textbf{Desglose de tareas}
\tarea{T14.1}{HU-14}{Comparar sílabo externo en movilidad}{D}{4 h}{13 nov. 2026}{13 nov. 2026}{Alex Jonathan Chisag}{Commit a9b5a8c: se incorporó la comparación del sílabo de otra universidad para movilidad externa.}
\textbf{Lo que queda funcionando}
\begin{itemize}\item El trámite externo acepta sílabos por materia y los usa en la comparación.\item Si no se adjunta uno, la materia queda para revisión docente sin comparación automática.\end{itemize}
\textbf{Pruebas de aceptación}
\begin{table}[H]\centering\small\begin{tabularx}{\textwidth}{|l|X|X|X|}\hline\rowcolor{grisTarjeta}\textbf{PA} & \textbf{Prueba} & \textbf{Entrada} & \textbf{Resultado esperado}\\\hline PA-14.1 & Comparación externa & Solicitud externa con sílabo & Compara unidades; sin documento informa que no hay sílabo para comparar.\\\hline\end{tabularx}\caption{Pruebas de aceptación de HU-14}\end{table}

\subsection{HU-15: Regla: créditos bajo 80~\% bloquean, contenidos bajo 80~\% solo recomiendan}
\historia{HU-15}{Docente}{Aplicar el mínimo de créditos y contenidos con efectos distintos}{Alta}{Alta}{8 SP}{14 al 25 de noviembre}{Alex Jonathan Chisag; Rolando Sebastián Miranda}{Como docente, quiero que los créditos bajo el mínimo bloqueen la homologación y que los contenidos bajos generen una recomendación, para aplicar la regla acordada con el docente.}{\item Dado que los créditos no llegan al 80~\%, cuando se analiza, entonces la materia queda bloqueada y no admite voto.\item Dado que solo los contenidos no llegan al 80~\%, cuando se analiza, entonces se recomienda no homologar y los docentes pueden decidir.}{La IA recomienda; los docentes conservan la decisión en los casos que no estén bloqueados por créditos.}{Aplicación de la regla del 80~\%}
\begin{figure}[H]\centering\begin{tikzpicture}[actor/.style={draw,ellipse},caso/.style={draw,ellipse,minimum width=3.4cm},node distance=1.1cm and 2cm]\node[actor](a){Docente};\node[caso,right=of a,yshift=.6cm](rev){Revisar porcentajes};\node[caso,below=of rev](bloq){Ver bloqueo por créditos};\node[caso,below=of bloq](rec){Ver recomendación por contenidos};\draw(a)--(rev);\draw(a)--(bloq);\draw(a)--(rec);\draw[dashed](rev)--node[right]{«extiende»}(bloq);\draw[dashed](rev)--node[right]{«extiende»}(rec);\end{tikzpicture}\caption{Casos de uso de HU-15}\end{figure}
\textbf{Desglose de tareas}
\tarea{T15.1}{HU-15}{Aplicar mínimos en la regla}{D}{6 h}{14 nov. 2026}{14 nov. 2026}{Alex Jonathan Chisag}{Commit 35d0796: se aplicó el mínimo del 80~\% en contenidos y créditos en la primera definición de la regla.}
\tarea{T15.2}{HU-15}{Separar recomendación y decisión docente}{M}{6 h}{24 nov. 2026}{24 nov. 2026}{Rolando Sebastián Miranda}{Commit 6aa9054: se dejó al agente recomendar y al docente decidir; también se incorporó el registro manual sin PDF.}
\tarea{T15.3}{HU-15}{Bloquear créditos insuficientes}{C}{4 h}{25 nov. 2026}{25 nov. 2026}{Alex Jonathan Chisag}{Commit 03eb265: se corrigió la regla para bloquear créditos bajos y dejar contenidos bajos como consejo.}
\tarea{T15.4}{HU-15}{Aclarar estados de recomendación}{C}{3 h}{25 nov. 2026}{25 nov. 2026}{Alex Jonathan Chisag}{Commit bcc90ec: se aclaró el resumen para distinguir las materias bloqueadas.}
\tarea{T15.5}{HU-15}{Corregir voto pendiente y créditos obligatorios}{C}{3 h}{25 nov. 2026}{25 nov. 2026}{Rolando Sebastián Miranda}{Commit cc3812d: se ajustó el aviso de decisión pendiente y el tratamiento de los créditos obligatorios.}
\tarea{T15.6}{HU-15}{Ajustar el texto de la regla en la base}{C}{2 h}{24 nov. 2026}{24 nov. 2026}{Rolando Sebastián Miranda}{Commit 17048b0: se corrigió el tamaño permitido para guardar la descripción de la regla RN-01.}
\textbf{Lo que queda funcionando}
\begin{itemize}\item Créditos por debajo del mínimo producen un bloqueo sin votación.\item Contenidos por debajo del mínimo producen una recomendación que los docentes pueden revisar y decidir.\item Profundidad y carga horaria se presentan como datos informativos.\end{itemize}
\textbf{Pruebas de aceptación}
\begin{table}[H]\centering\small\begin{tabularx}{\textwidth}{|l|X|X|X|}\hline\rowcolor{grisTarjeta}\textbf{PA} & \textbf{Prueba} & \textbf{Entrada} & \textbf{Resultado esperado}\\\hline PA-15.1 & Créditos inferiores al 80~\% & Materia con créditos insuficientes & Queda bloqueada y sin controles de voto.\\\hline PA-15.2 & Solo contenidos inferiores al 80~\% & Créditos suficientes y baja cobertura & Se muestra recomendación y se permiten los votos docentes.\\\hline\end{tabularx}\caption{Pruebas de aceptación de HU-15}\end{table}

\subsection{HU-16: Decisión de los dos docentes y decisión final del administrador cuando no coinciden}
\historia{HU-16}{Docentes y Administrador}{Registrar las dos decisiones docentes y resolver desacuerdos}{Alta}{Alta}{8 SP}{13 al 25 de noviembre}{Alex Jonathan Chisag; Rolando Sebastián Miranda}{Como docente revisor, quiero registrar mi voto por materia; como administrador, quiero resolver con un motivo los desacuerdos, para que cada solicitud alcance una decisión final trazable.}{\item Dada una materia revisable, cuando vota cada docente, entonces queda guardada su decisión.\item Dados votos distintos, cuando se presenta el resultado, entonces el administrador puede decidir y debe explicar el motivo.\item Dada una decisión final, cuando se intenta cambiar un voto, entonces el sistema lo impide.}{Los votos docentes se conservan como parte del expediente.}{Votos y decisión final}
\begin{figure}[H]\centering\begin{tikzpicture}[actor/.style={draw,ellipse},caso/.style={draw,ellipse,minimum width=3.6cm},node distance=1.1cm and 2.1cm]\node[actor](d){Docente};\node[actor,below=2.1cm of d](a){Administrador};\node[caso,right=of d](voto){Registrar voto};\node[caso,right=of a](final){Resolver desacuerdo};\node[caso,below=of voto](estado){Consultar resultado};\draw(d)--(voto);\draw(a)--(final);\draw(d)--(estado);\draw(a)--(estado);\draw[dashed](final)--node[right]{«extiende»}(estado);\end{tikzpicture}\caption{Casos de uso de HU-16}\end{figure}
\textbf{Desglose de tareas}
\tarea{T16.1}{HU-16}{Avisar sobre sesión y solicitudes docentes}{M}{4 h}{13 nov. 2026}{13 nov. 2026}{Alex Jonathan Chisag}{Commit 68a2e56: se agregó aviso de sesión próxima a vencer y aviso de nuevas solicitudes al docente.}
\tarea{T16.2}{HU-16}{Añadir decisión final del administrador}{D}{6 h}{24 nov. 2026}{24 nov. 2026}{Rolando Sebastián Miranda}{Commit 059c63a: se implementó la decisión final administrativa cuando los docentes no coinciden y el reenvío al agente.}
\tarea{T16.3}{HU-16}{Cerrar votos tras decisión final}{C}{4 h}{25 nov. 2026}{25 nov. 2026}{Alex Jonathan Chisag}{Commit 1d8798c: se cerraron los votos después de la decisión final y se aclaró el estado.}
\textbf{Lo que queda funcionando}
\begin{itemize}\item Cada docente registra su voto y el desacuerdo queda visible.\item El administrador resuelve el desacuerdo con un motivo.\item Una vez resuelto, los votos no pueden modificarse y la solicitud muestra su estado final.\end{itemize}
\textbf{Pruebas de aceptación}
\begin{table}[H]\centering\small\begin{tabularx}{\textwidth}{|l|X|X|X|}\hline\rowcolor{grisTarjeta}\textbf{PA} & \textbf{Prueba} & \textbf{Entrada} & \textbf{Resultado esperado}\\\hline PA-16.1 & Voto docente & Decisión del primer revisor & El voto queda guardado.\\\hline PA-16.2 & Sin acuerdo & Dos votos diferentes & La materia queda pendiente de decisión administrativa.\\\hline PA-16.3 & Decisión final con motivo & Decisión administrativa y motivo & Se guarda la decisión final y el motivo.\\\hline PA-16.4 & Votos cerrados & Intento de cambiar voto tras decisión & El sistema impide el cambio.\\\hline\end{tabularx}\caption{Pruebas de aceptación de HU-16}\end{table}

\subsection{HU-17: Tabla RESULTADO\_HOMOLOGACION con el resultado por materia}
\historia{HU-17}{Administrador}{Consultar el resultado de homologación por materia}{Media}{Media}{5 SP}{13 de noviembre}{Alex Jonathan Chisag}{Como administrador, quiero consultar una fila por materia con los porcentajes, la recomendación y el resultado final, para revisar los datos del proceso de forma ordenada.}{\item Dado un análisis terminado, cuando se consulta la tabla, entonces existe una fila por cada materia solicitada.\item Dada una decisión docente o administrativa, cuando cambia el resultado, entonces la fila refleja el resultado actualizado.}{La tabla complementa el análisis completo de la solicitud.}{Resultado por materia}
\begin{figure}[H]\centering\begin{tikzpicture}[actor/.style={draw,ellipse},caso/.style={draw,ellipse,minimum width=3.4cm},node distance=1.1cm and 2cm]\node[actor](a){Administrador};\node[caso,right=of a,yshift=.6cm](consulta){Consultar resultados};\node[caso,below=of consulta](detalle){Ver resultado por materia};\draw(a)--(consulta);\draw(a)--(detalle);\draw[dashed](consulta)--node[right]{«incluye»}(detalle);\end{tikzpicture}\caption{Casos de uso de HU-17}\end{figure}
\textbf{Desglose de tareas}
\tarea{T17.1}{HU-17}{Guardar resultados por materia}{D}{4 h}{13 nov. 2026}{13 nov. 2026}{Alex Jonathan Chisag}{Commit 5596b30: se creó la tabla RESULTADO\_HOMOLOGACION y se guardó el mínimo del reglamento en el resultado.}
\textbf{Lo que queda funcionando}
\begin{itemize}\item Hay una fila de resultado por materia, con la recomendación, los porcentajes y las decisiones.\item El resultado se actualiza conforme se completan las revisiones.\end{itemize}
\textbf{Pruebas de aceptación}
\begin{table}[H]\centering\small\begin{tabularx}{\textwidth}{|l|X|X|X|}\hline\rowcolor{grisTarjeta}\textbf{PA} & \textbf{Prueba} & \textbf{Entrada} & \textbf{Resultado esperado}\\\hline PA-17.1 & Resultado por materia & Análisis con varias materias & Existe una fila por materia con porcentajes, votos y resultado final.\\\hline\end{tabularx}\caption{Pruebas de aceptación de HU-17}\end{table}

\subsection{HU-18: Registro manual sin PDF, lo más automático posible y solo con materias que se relacionan}
\historia{HU-18}{Operador}{Registrar manualmente una solicitud y completar los datos posibles}{Media}{Alta}{5 SP}{24 al 25 de noviembre}{Rolando Sebastián Miranda; Jair León}{Como operador, quiero registrar los datos del certificado sin PDF y recibir sugerencias relacionadas, para atender casos en que no dispongo del archivo y reducir la escritura manual.}{\item Dada una carrera y un nivel, cuando se carga la malla, entonces se completa el historial con sus materias.\item Dada una materia aprobada sin relación con el destino, cuando se revisa la solicitud, entonces se muestra un aviso y no se registra esa relación.}{El registro manual utiliza controles equivalentes a los del certificado.}{Registro manual}
\begin{figure}[H]\centering\begin{tikzpicture}[actor/.style={draw,ellipse},caso/.style={draw,ellipse,minimum width=3.5cm},node distance=1.1cm and 2cm]\node[actor](a){Operador};\node[caso,right=of a,yshift=.6cm](datos){Ingresar datos};\node[caso,below=of datos](carga){Cargar materias de malla};\node[caso,below=of carga](valida){Validar relación de materias};\draw(a)--(datos);\draw(a)--(carga);\draw(a)--(valida);\draw[dashed](datos)--node[right]{«incluye»}(valida);\end{tikzpicture}\caption{Casos de uso de HU-18}\end{figure}
\textbf{Desglose de tareas}
\tarea{T18.1}{HU-18}{Mantener número revisado en registro manual}{C}{3 h}{24 nov. 2026}{24 nov. 2026}{Rolando Sebastián Miranda}{Commit 37bca1f: se corrigió para que el número revisado sea el mismo que queda registrado.}
\tarea{T18.2}{HU-18}{Completar registro manual y sugerencias relacionadas}{M}{4 h}{25 nov. 2026}{25 nov. 2026}{Rolando Sebastián Miranda}{Commit 28e15e5: se añadió el llenado automático y se limitaron las solicitudes a materias relacionadas.}
\tarea{T18.3}{HU-18}{Listas para nivel, materia por nivel y períodos de la UTA}{M}{4 h}{25 nov. 2026}{25 nov. 2026}{Jair León}{Commit ac93ec4: el nivel y la materia se eligen de listas (la materia se filtra por nivel) y los períodos académicos de la UTA vienen precargados.}
\textbf{Lo que queda funcionando}
\begin{itemize}\item Se puede ingresar una solicitud sin adjuntar el PDF y completar el historial desde la malla.\item Las sugerencias ayudan a seleccionar materias relacionadas y se informan las que no guardan relación.\end{itemize}
\textbf{Pruebas de aceptación}
\begin{table}[H]\centering\small\begin{tabularx}{\textwidth}{|l|X|X|X|}\hline\rowcolor{grisTarjeta}\textbf{PA} & \textbf{Prueba} & \textbf{Entrada} & \textbf{Resultado esperado}\\\hline PA-18.1 & Llenado automático & Carrera de origen y nivel & Se cargan materias de la malla y sus datos disponibles.\\\hline PA-18.2 & Materias no relacionadas & Materia de origen sin afinidad al destino & Se explica el problema y no se registra esa relación.\\\hline\end{tabularx}\caption{Pruebas de aceptación de HU-18}\end{table}

\subsection{HU-19: Cuentas y sesión: restablecer contraseña, aviso de sesión por vencer y solo entrar a páginas permitidas}
\historia{HU-19}{Administrador}{Administrar cuentas y proteger el acceso por rol}{Media}{Media}{5 SP}{13 al 22 de noviembre}{Alex Jonathan Chisag}{Como administrador, quiero restablecer contraseñas y como usuario quiero recibir un aviso de sesión próxima a vencer y volver solo a páginas permitidas, para operar con acceso controlado.}{\item Dada una cuenta, cuando el administrador restablece su contraseña, entonces puede definir una nueva.\item Dada una sesión próxima a vencer, cuando el usuario está activo, entonces recibe un aviso.\item Dado un destino no permitido para el rol, cuando inicia sesión, entonces el sistema lo redirige a una página autorizada.}{La autorización depende del rol asignado.}{Cuentas y sesión}
\begin{figure}[H]\centering\begin{tikzpicture}[actor/.style={draw,ellipse},caso/.style={draw,ellipse,minimum width=3.5cm},node distance=1.1cm and 2cm]\node[actor](a){Administrador};\node[actor,below=1.5cm of a](u){Usuario};\node[caso,right=of a](clave){Restablecer contraseña};\node[caso,right=of u](sesion){Avisar sesión próxima a vencer};\node[caso,below=of sesion](acceso){Validar página por rol};\draw(a)--(clave);\draw(u)--(sesion);\draw(u)--(acceso);\end{tikzpicture}\caption{Casos de uso de HU-19}\end{figure}
\textbf{Desglose de tareas}
\tarea{T19.1}{HU-19}{Gestionar contraseñas y avisos de sesión}{M}{4 h}{22 nov. 2026}{22 nov. 2026}{Alex Jonathan Chisag}{Commit 0a3f102: se incorporaron avisos claros, gestión de contraseñas y ajustes del menú.}
\textbf{Lo que queda funcionando}
\begin{itemize}\item El administrador dispone de la acción de restablecer contraseña.\item El usuario recibe un aviso de sesión próxima a vencer y el retorno tras iniciar sesión considera su rol.\end{itemize}
\textbf{Pruebas de aceptación}
\begin{table}[H]\centering\small\begin{tabularx}{\textwidth}{|l|X|X|X|}\hline\rowcolor{grisTarjeta}\textbf{PA} & \textbf{Prueba} & \textbf{Entrada} & \textbf{Resultado esperado}\\\hline PA-19.1 & Restablecer contraseña & Cuenta seleccionada por administrador & Permite definir una contraseña nueva.\\\hline PA-19.2 & Página permitida & Inicio de sesión con rol definido & Solo redirige a una ruta autorizada para ese rol.\\\hline\end{tabularx}\caption{Pruebas de aceptación de HU-19}\end{table}
